\begin{figure}[tb]
\centering
\begin{tikzpicture}
\begin{axis}[axis main,width=.78\linewidth,height=6cm,
 xlabel={torque $M$},ylabel={minimum loss $P_{\min}$},xmin=0,xmax=5.2,ymin=0,ymax=5.2]
 \addplot[loss curve,domain=0:4.8,samples=100]{.34*x+.12*x^2};
 \addplot[support line] coordinates {(0,3.25) (5.1,3.25)};
 \addplot[support line] coordinates {(3.72,0) (3.72,5)};
 \addplot[only marks,mark=*,mark size=2.2pt,ForestGreen] coordinates {(2.45,1.55)};
 \addplot[only marks,mark=*,mark size=2.2pt,BrickRed] coordinates {(3.72,3.25)};
 \node[annotation] at (axis cs:1.45,1.95){tracking};
 \node[annotation] at (axis cs:4.3,3.75){maximum feasible};
 \node[annotation] at (axis cs:1.2,3.55){$P_{\mathrm{Cu},\max}$};
\end{axis}
\end{tikzpicture}
\caption{Minimum loss and maximum torque are inverse readings of one strictly
increasing efficient frontier on a fixed torque-sign branch.}
\label{fig:frontier}
\end{figure}
